\section{Flood Fill}
\label{sec:graph-flood-fill}

\problemstatement{Given a 2D grid of pixel colors, a starting pixel
$(\textit{sr},\textit{sc})$, and a $\textit{newColor}$, change the
color of the starting pixel and every pixel \textbf{4-directionally
connected} to it through pixels of the \textbf{same original color} to
$\textit{newColor}$.}

\begin{patternbox}
\emph{``Change every connected pixel of the same color''} is a direct
single-source DFS (or BFS) on the implicit grid graph, where two
pixels are connected exactly when they are $4$-directional neighbors
\textbf{and} share the starting pixel's original color.
\end{patternbox}

\subsection*{Algorithm}

\begin{enumerate}
  \item Record $\textit{originalColor} = \textit{image}[\textit{sr}][\textit{sc}]$.
  \item If $\textit{originalColor} = \textit{newColor}$, return the
        image unchanged immediately (otherwise the DFS below would
        recurse infinitely, since every newly-colored pixel would
        still match the ``original color'' check).
  \item DFS from $(\textit{sr},\textit{sc})$: if the current pixel is
        in bounds and equals $\textit{originalColor}$, set it to
        $\textit{newColor}$, then recurse into all $4$ neighbors.
\end{enumerate}

\subsection*{C++ implementation}

\begin{lstlisting}
#include <bits/stdc++.h>
using namespace std;

void fill(vector<vector<int>>& image, int r, int c, int originalColor, int newColor) {
    int rows = image.size(), cols = image[0].size();
    if (r < 0 || r >= rows || c < 0 || c >= cols) return;
    if (image[r][c] != originalColor) return;

    image[r][c] = newColor;
    fill(image, r - 1, c, originalColor, newColor);
    fill(image, r + 1, c, originalColor, newColor);
    fill(image, r, c - 1, originalColor, newColor);
    fill(image, r, c + 1, originalColor, newColor);
}

vector<vector<int>> floodFill(vector<vector<int>>& image, int sr, int sc, int newColor) {
    int originalColor = image[sr][sc];
    if (originalColor != newColor) {
        fill(image, sr, sc, originalColor, newColor);
    }
    return image;
}

int main() {
    vector<vector<int>> image = {
        {1, 1, 1},
        {1, 1, 0},
        {1, 0, 1}
    };
    vector<vector<int>> result = floodFill(image, 1, 1, 2);
    for (auto& row : result) {
        for (int x : row) cout << x << " ";
        cout << endl;
    }
    // 2 2 2
    // 2 2 0
    // 2 0 1
    return 0;
}
\end{lstlisting}

\subsection*{Dry run}

\dryrun{Take the $3\times3$ image above, starting at $(1,1)$ (value
$1$), $\textit{newColor}=2$.}

$\textit{originalColor}=1 \neq 2$, proceed. Fill$(1,1)$: matches,
set to $2$; recurse to $(0,1)$: matches, set $2$; recurse further to
$(0,0)$ and $(0,2)$, both matching, set $2$. Back to $(1,1)$'s other
neighbors: $(2,1)=0$, no match, skip. $(1,0)=1$, matches, set $2$;
its neighbor $(2,0)=1$, matches, set $2$. $(1,2)=0$, no match, skip.
Final grid: top-left $2\times2$ block plus $(0,2)$ and $(2,0)$ all
become $2$, while the two $0$-pixels and the isolated $(2,2)=1$ pixel
(not $4$-connected to the start through same-colored pixels) remain
unchanged.

\begin{complexitybox}
\textbf{Time Complexity.} $O(\textit{rows} \times \textit{cols})$ --
each pixel is visited (and re-colored) at most once, since after being
set to $\textit{newColor}$ it no longer matches $\textit{originalColor}$
and further recursive calls on it return immediately.
\textbf{Space Complexity.} $O(\textit{rows} \times \textit{cols})$
recursion stack in the worst case (an entirely same-colored grid).
\end{complexitybox}

\begin{takeawaybox}
Flood Fill is the simplest possible instance of ``traverse an implicit
grid graph'' -- a single source, a single connectivity condition (same
original color), and no auxiliary bookkeeping beyond the grid itself
(the grid doubles as its own \textit{visited} array, since a re-colored
pixel naturally fails the match check on revisit).
\end{takeawaybox}
